\begin{titlepage}
    \begin{center}
        \vspace*{1cm}

        \huge
        \textbf{Calibration and Image Reconstruction Techniques
            for FMCW MIMO Radars with Consideration of Antenna Radiation Patterns}
        \vspace{0.5cm}

        \large
        Kalibrations- und Bildrekonstruktionsverfahren für FMCW-MIMO-Radare unter Berücksichtigung der Antennenabstrahlungscharakteristik
        \vspace{1.5cm}

        \textbf{Daniel Gotzens}

        \vfill

        \normalsize
        Masterarbeit vorgelegt dem \\

        \vspace{0.8cm}

        \includegraphics[width=4cm]{../figures/IHF-Logo_platzhalter.png}

        Lehrstuhl für Radarsystemtechnik \\
        Institut für Hochfrequenztechnik \\
        RWTH Aachen University\\
        \today

        \vspace{1cm}
        In Kooperation mit \\

        \includegraphics[height=18pt]{../figures/indurad-logo.png}

        indurad GmbH
        Aachen

    \end{center}
\end{titlepage}

\begin{abstract}
    Increased image quality can be achieved by implementing backprojection on a millimeter-wave TDM FMCW MIMO Radar.
    In this thesis, a backprojection-based imaging algorithm is proposed and implemented for a radar sensor.
    This radar sensor's functionality is tested, the sensor is calibrated,
    and then the algorithm is compared to the current FFT-based algorithm as well as MUSIC.
    It is found that near-field accuracy of the image generated with backprojection and MUSIC is increased,
    and that real-time capability can only be guaranteed for backprojection and the current FFT-based approach.
\end{abstract}